The competing families cover graph traversal (CAGRA fp32/int8),
product-quantised IVF and GPU IVF-RaBitQ \cite{cagra,pq,ivfrabitq}.
FP32 GPU brute force provides an exhaustive-search reference. Each family
sweeps its own implementation parameters (\cref{tab:3-grids}); the recorded
index configurations can differ in storage and reconstruction fidelity.
Consequently the main frontier measures achievable speed at a quality
target within these grids. A separate memory account accompanies it.

\par\medskip\noindent\begin{minipage}{\linewidth}
\centering\captionsetup{hypcap=false}
\captionof{table}[Principal E1 tuning dimensions]{Principal E1 tuning dimensions. Exact per-point parameters and
coverage extensions remain in the raw records. The grid is finite and
differs between methods.}
\label{tab:3-grids}
\footnotesize\setlength{\tabcolsep}{4pt}
\begin{tabularx}{\linewidth}{@{}>{\raggedright\arraybackslash}p{.22\linewidth}>{\raggedright\arraybackslash}X@{}}
\toprule
Family & Recorded search dimensions \\
\midrule
JHQ & Probes and $\alpha$; initial probes 64/128/256 and
$0.5/1/2$ times the workload budget; coverage extensions reach probe
1,024 and larger/smaller budgets \\
CAGRA fp32/int8 & Graph degree 32/64; search width 1/2/4/8;
intermediate top-$k$ capacity 32--512 \\
IVF-PQ & Eight-bit subcodes; 192/384/768-byte codes at $d=768$,
384/768/1,536 at higher dimensions; fp16/fp32 lookup settings;
probes 16--512 \\
IVF-RaBitQ & 4/8/9 bits per dimension, supported search modes,
probes 16--1,024; additional 4,096-list OpenAI indexes \\
GPU brute force & Exhaustive FP32 squared-L2 search, common host boundary \\
\bottomrule
\end{tabularx}
\end{minipage}\par\medskip

The frontier retains points not dominated in both recall and QPS: point
$a$ dominates $b$ when $R_a\geq R_b$ and $Q_a\geq Q_b$, with at least one
strict inequality. This envelope can combine different index configurations.
For target recall $r$, the measured estimator is
\begin{equation}
Q_{\rm meas}(r)=\max\{Q_i:R_i\geq r\}.
\label{eq:3-measured-frontier}
\end{equation}
It is undefined when no sampled point qualifies. The interpolated estimator
uses adjacent frontier points $(R_a,Q_a)$ and $(R_b,Q_b)$:
\begin{equation}
\log Q_{\rm int}(r)=(1-\lambda)\log Q_a+\lambda\log Q_b,
\qquad \lambda=\frac{r-R_a}{R_b-R_a}.
\label{eq:3-interpolated-frontier}
\end{equation}
The archived parser first accepts a frontier point within $5\times10^{-4}$
of the target as a direct match, including near endpoints; otherwise it
interpolates only inside the measured range. It extrapolates no segment.
Both estimators are reported where comparisons are close: one selects an
observed qualifying configuration, while the other reduces sensitivity to
unequal grid spacing. Neither supplies an equal-storage comparison.
T5 also retains the fastest qualifying measured endpoint when a target
lies below a method's minimum frontier recall; interpolation being
undefined does not make a more accurate configuration ineligible.

The parser retains failures and invalid-output diagnostics outside the
valid-point frontier. Repeat-call ID changes alone do not discard a point;
set changes and ordering changes are distinguished in the source records.
Corrected isolated large-baseline reruns supersede contaminated builds,
and Stella runs skipped at batch 1,000 contribute no search points.
\Cref{sec:a-e1-protocol} records those decisions. The thesis reconstructs
its frontier from all declared baseline families, including int8 CAGRA and
adverse operating regions (\cref{fig:7-frontiers}). Historical ADC and poster
plots retain their own version and measurement boundaries.
